
\subsection{含幺情形 I：无滤过版本}
\label{subsec:4-4-4}

在转向滤过版本之前，先说明如何修改上面的论证以处理含幺情形。假设
$\overline{\mathfrak f}$ 含幺。我们将构造
$\overline{\mathfrak F}$，使它也含幺，并且在
$\overline{\mathfrak f}$ 之上与 $\overline{\mathrm{Incl}}$、
$\overline{\mathrm{Eval}}$ 相容。令 $\overline C$ 为以 $e$ 为幺元的
$A_\infty$ 代数。

\begin{definition}{4.4.22}\label{def:4-4-22}
置 $\overline C_{\mathrm{red}}=\overline C/R[e]$，并定义如下的
\emph{部分约化 bar 复形}\liftIndex{I162}：
\begin{equation}
 B_k^{\mathrm{red}}\overline C[1]
 =
 \begin{cases}
  \overline C[1]\otimes
  B_{k-2}\overline C_{\mathrm{red}}[1]\otimes\overline C[1],
  &k\geq2,\\
  \overline C[1],&k=1,
 \end{cases}
\tag{4.4.23}\label{eq:4-4-23}
\end{equation}
以及
\[
 \widehat B^{\mathrm{red}}\overline C[1]
 =\widehat{\bigoplus}_{k\geq1}
 B_k^{\mathrm{red}}\overline C[1].
\]
\end{definition}

\begin{statement}{引理--定义}{4.4.24}\label{lem-def:4-4-24}
$\widehat{\overline d}\colon
B\overline C[1]\to B\overline C[1]$ 在
$\widehat B^{\mathrm{red}}\overline C[1]$ 上诱导微分
$\widehat{\overline d}^{\mathrm{red}}$。
\end{statement}

\begin{proof}
考虑元素
$\mathbf x_1\otimes e\otimes\mathbf x_2\in B\overline C[1]$，其中
$\mathbf x_i\in B_{k_i}\overline C[1]$、$k_i\geq1$。写
$\mathbf x_1=\mathbf x_{1-}\otimes x_{\mathrm{last}}$、
$\mathbf x_2=x_{\mathrm{first}}\otimes\mathbf x_{2-}$。于是
\[
 \begin{aligned}
 \widehat{\overline d}
 (\mathbf x_1\otimes e\otimes\mathbf x_2)
 &={}
 (-1)^{\deg'\mathbf x_{1-}}
 \mathbf x_{1-}\otimes
 \overline{\mathfrak m}_2(x_{\mathrm{last}},e)
 \otimes\mathbf x_2\\
 &\quad{}+
 (-1)^{\deg'\mathbf x_1}
 \mathbf x_1\otimes
 \overline{\mathfrak m}_2(e,x_{\mathrm{first}})
 \otimes\mathbf x_{2-}=0
 \end{aligned}
\]
于 $\widehat B^{\mathrm{red}}\overline C[1]$ 中；这里用了条件 3.2.18。
引理得证。
\translatornote{原书在第二项的次数指数和张量因子中都写了未定义的
$\mathbf x$；由本段刚给出的分解以及相邻第一项可知，两处均应为
$\mathbf x_1$，译文据此补回下标。}
\end{proof}

注意，
$\widehat{\overline d}\colon
\widehat B\overline C_{\mathrm{red}}[1]
\to\widehat B\overline C_{\mathrm{red}}[1]$
并没有良定义。

下面定义算子
\[
 s_{\mathrm{left}},s_{\mathrm{right}}\colon
 B_k^{\mathrm{red}}\overline C[1]
 \longrightarrow B_{k+1}^{\mathrm{red}}\overline C[1]
\]
如下：
\[
 s_{\mathrm{left}}(\mathbf x)=e\otimes\mathbf x,
 \qquad
 s_{\mathrm{right}}(\mathbf x)
 =(-1)^{\deg\mathbf x}\mathbf x\otimes e.
\]

\begin{lemma}{4.4.25}\label{lem:4-4-25}
\[
 \widehat{\overline d}^{\mathrm{red}}\circ s_{\mathrm{left}}
 +s_{\mathrm{left}}\circ\widehat{\overline d}^{\mathrm{red}}
 =
 \widehat{\overline d}^{\mathrm{red}}\circ s_{\mathrm{right}}
 +s_{\mathrm{right}}\circ\widehat{\overline d}^{\mathrm{red}}
 =\operatorname{id}.
\]
\end{lemma}

证明只需利用条件 3.2.18 作直接计算。

令 $\overline C_1,\overline C_2$ 为含幺 $A_\infty$ 代数。假设有次数为
$0$ 的一列同态
\[
 \overline{\mathfrak f}_{k,\mathrm{red}}\colon
 B_k^{\mathrm{red}}\overline C_1[1]
 \longrightarrow\overline C_2[1],
 \qquad
 \overline{\mathfrak f}_{1,\mathrm{red}}(e)=e.
\]
它们按如下方式诱导映射
\begin{equation}
 \widehat{\overline{\mathfrak f}}_{\mathrm{red}}\colon
 \widehat B^{\mathrm{red}}\overline C_1[1]
 \longrightarrow\widehat B\overline C_2[1].
\tag{4.4.26}\label{eq:4-4-26}
\end{equation}
注意，右端是通常的 bar 复形，而不是约化 bar 复形。选取分裂
\begin{equation}
 \overline C=\overline C_{\mathrm{red}}\oplus R[e].
\tag{4.4.27}\label{eq:4-4-27}
\end{equation}
参见定义 3.2.20 (3)。利用 \eqref{eq:4-4-27}，把
$B_k\overline C_{1,\mathrm{red}}[1]$ 看成
$\widehat B_k^{\mathrm{red}}\overline C_1[1]$ 的子模；先在
$\widehat B\overline C_{1,\mathrm{red}}[1]$ 上，用与 (3.2.28) 相同的
公式定义 $\widehat{\overline{\mathfrak f}}_{\mathrm{red}}$，再置
\translatornote{原书此句把正在构造的
$\widehat{\overline{\mathfrak f}}_{\mathrm{red}}$ 简写为
$\widehat{\overline{\mathfrak f}}$，并在定义域中漏排移位 $[1]$；
式 (4.4.26) 及紧随其后的四个延拓式唯一确定了这里的记号，译文统一
补回。}
\[
 \begin{aligned}
 \widehat{\overline{\mathfrak f}}_{\mathrm{red}}
 (e\otimes\mathbf x)
 &=e\otimes
 \widehat{\overline{\mathfrak f}}_{\mathrm{red}}(\mathbf x),&
 \widehat{\overline{\mathfrak f}}_{\mathrm{red}}
 (\mathbf x\otimes e)
 &=\widehat{\overline{\mathfrak f}}_{\mathrm{red}}(\mathbf x)
 \otimes e,\\
 \widehat{\overline{\mathfrak f}}_{\mathrm{red}}
 (e\otimes\mathbf x\otimes e)
 &=e\otimes
 \widehat{\overline{\mathfrak f}}_{\mathrm{red}}(\mathbf x)
 \otimes e,&
 \widehat{\overline{\mathfrak f}}_{\mathrm{red}}(e)&=e,
 \end{aligned}
\]
其中 $\mathbf x\in\widehat B\overline C_{1,\mathrm{red}}[1]$。又定义
$\widehat{\overline{\mathfrak f}}\colon
\widehat B\overline C_1[1]\to\widehat B\overline C_2[1]$：
\begin{equation}
 \begin{aligned}
 &\widehat{\overline{\mathfrak f}}
 (\mathbf x_1\otimes e\otimes\cdots\otimes e\otimes\mathbf x_\ell)\\
 &\qquad=
 \widehat{\overline{\mathfrak f}}_{\mathrm{red}}(\mathbf x_1)
 \otimes e\otimes\cdots\otimes e\otimes
 \widehat{\overline{\mathfrak f}}_{\mathrm{red}}(\mathbf x_\ell),
 \end{aligned}
\tag{4.4.28}\label{eq:4-4-28}
\end{equation}
其中
$\mathbf x_i\in B_{k_i}\overline C_{1,\mathrm{red}}[1]$。仍用
\eqref{eq:4-4-27} 把
$B_k\overline C_{1,\mathrm{red}}[1]$ 看成
$\widehat B_k^{\mathrm{red}}\overline C_1[1]$ 的子模，并对不属于
$\widehat B^{\mathrm{red}}\overline C_1[1]$ 的
$1\in B_0\overline C_{1,\mathrm{red}}[1]$ 置
$\widehat{\overline{\mathfrak f}}_{\mathrm{red}}(1)=1$。

\begin{remark}{4.4.29}\label{rem:4-4-29}
一般而言，$\overline C_{\mathrm{red}}$ 不是 $A_\infty$ 子代数。事实上，
$\overline{\mathfrak m}_k$ 在
$B_k\overline C_{\mathrm{red}}[1]$ 上的限制之像可能为 $ce$。而且在纯
代数情形下，分裂 \eqref{eq:4-4-27} 未必是典范的。（在本书的几何
情形下存在典范分裂，但这里不使用它。）
\end{remark}

\begin{lemma}{4.4.30}\label{lem:4-4-30}
令 $\overline{\mathfrak f}_{k,\mathrm{red}}$ 为式
\eqref{eq:4-4-26} 中的一列同态。假设
\begin{equation}
 \widehat{\overline{\mathfrak f}}_{\mathrm{red}}
 \circ\widehat{\overline d}^{\mathrm{red}}
 =\widehat{\overline d}
 \circ\widehat{\overline{\mathfrak f}}_{\mathrm{red}}.
\tag{4.4.31.1}\label{eq:4-4-31-1}
\end{equation}
则式 \eqref{eq:4-4-28} 给出的同态是含幺 $A_\infty$ 同态。反之，每个
含幺 $A_\infty$ 同态都来自一列满足 \eqref{eq:4-4-31-1} 以及
\begin{equation}
 \overline{\mathfrak f}_{k,\mathrm{red}}\circ s_{\mathrm{left}}
 =\overline{\mathfrak f}_{k,\mathrm{red}}\circ s_{\mathrm{right}}=0
\tag{4.4.31.2}\label{eq:4-4-31-2}
\end{equation}
的 $\overline{\mathfrak f}_{k,\mathrm{red}}$。
\end{lemma}

\begin{proof}
假设 (4.4.31) 成立。对
$\mathbf x_i\in B_{k_i}\overline C_{1,\mathrm{red}}[1]$，计算
\begin{equation}
 \begin{aligned}
 &\overline{\mathfrak m}\bigl(
 \widehat{\overline{\mathfrak f}}
 (\mathbf x_1\otimes e\otimes\cdots\otimes e\otimes\mathbf x_\ell)
 \bigr)\\
 &\qquad=
 \overline{\mathfrak m}\bigl(
 \widehat{\overline{\mathfrak f}}_{\mathrm{red}}(\mathbf x_1)
 \otimes e\otimes\cdots\otimes e\otimes
 \widehat{\overline{\mathfrak f}}_{\mathrm{red}}(\mathbf x_\ell)
 \bigr).
 \end{aligned}
\tag{4.4.32}\label{eq:4-4-32}
\end{equation}

情形 1：若 $k_1,k_\ell>0$，则 \eqref{eq:4-4-32} 为零；同时
\begin{equation}
 (\overline{\mathfrak f}\circ\widehat{\overline d})
 (\mathbf x_1\otimes e\otimes\cdots\otimes e\otimes\mathbf x_\ell)=0.
\tag{4.4.33}\label{eq:4-4-33}
\end{equation}

情形 2：若 $\ell>2$，则 \eqref{eq:4-4-32} 为零，而
\eqref{eq:4-4-33} 也成立。

情形 3：若 $\ell=1$，则
\[
 (\widehat{\overline d}\circ
 \widehat{\overline{\mathfrak f}})(\mathbf x_1)
 =
 (\widehat{\overline d}\circ
 \widehat{\overline{\mathfrak f}}_{\mathrm{red}})(\mathbf x_1)
 =
 (\widehat{\overline{\mathfrak f}}_{\mathrm{red}}
 \circ\widehat{\overline d}^{\mathrm{red}})(\mathbf x_1).
\]
取其分量，得到
\[
 (\overline{\mathfrak m}\circ
 \widehat{\overline{\mathfrak f}})(\mathbf x_1)
 =
 (\overline{\mathfrak f}\circ\widehat{\overline d})(\mathbf x_1),
\]
因为
$\widehat{\overline{\mathfrak f}}_{\mathrm{red}}$ 的
$B_1\overline C_2[1]$-分量就是 $\overline{\mathfrak f}$。

情形 4：若 $\ell=2$、$k_1=0$，则
\[
 \begin{aligned}
 \overline{\mathfrak m}
 \bigl(\widehat{\overline{\mathfrak f}}(e\otimes\mathbf x)\bigr)
 &=\overline{\mathfrak m}
 \bigl(e\otimes
 \widehat{\overline{\mathfrak f}}_{\mathrm{red}}(\mathbf x)\bigr)\\
 &=\overline{\mathfrak m}_2
 \bigl(e\otimes\overline{\mathfrak f}(\mathbf x)\bigr)
 =\overline{\mathfrak f}(\mathbf x).
 \end{aligned}
\]
另一方面，由
$\widehat{\overline d}(e\otimes\mathbf x)
=\mathbf x-e\otimes\widehat{\overline d}(\mathbf x)$，有
\[
 \overline{\mathfrak f}
 \bigl(\widehat{\overline d}(e\otimes\mathbf x)\bigr)
 =\overline{\mathfrak f}(\mathbf x).
\]

情形 5：$\ell=2$、$k_2=0$ 的情形与情形 4 类似。

至此证明了 $\widehat{\overline{\mathfrak f}}$ 是上链同态。另一方面，
\eqref{eq:4-4-28} 蕴含它是余代数同态，因而
$\overline{\mathfrak f}$ 是 $A_\infty$ 同态；容易看出它含幺。逆命题
也很容易。
\translatornote{原书在这里把刚构造的同态称为“滤过
$A_\infty$ 同态”；本小节明确处理 $R$ 上的无滤过代数，且直到
4.4.5 才转向能量滤过，故译文删去“滤过”。}
\end{proof}

现在可以证明定理 4.2.34 的无滤过含幺版本了。令
$\overline{\mathfrak f}_{k,\mathrm{red}}$ 如引理 4.4.30，并令
$\overline{\mathfrak C}_1,\overline{\mathfrak C}_2$ 分别为
$[0,1]\times\overline C_1$、$[0,1]\times\overline C_2$ 的模型。选取与
$\overline{\mathrm{Incl}}$、$\overline{\mathrm{Eval}}$ 相容的分裂
\eqref{eq:4-4-27}。

\begin{proposition}{4.4.34}\label{prop:4-4-34}
对 $k=1,\ldots,K$，存在
\[
 \overline{\mathfrak F}_{k,\mathrm{red}}\colon
 B_k^{\mathrm{red}}\overline{\mathfrak C}_1[1]
 \longrightarrow\overline{\mathfrak C}_2[1]
\]
满足 (4.4.31)，使得所诱导的含幺 $A_K$ 同态
$\overline{\mathfrak F}^{(K)}\colon
\overline{\mathfrak C}_1\to\overline{\mathfrak C}_2$，在
$\overline{\mathfrak f}$ 之上与 $\overline{\mathrm{Incl}}$、
$\overline{\mathrm{Eval}}$ 相容。
\translatornote{原书把本命题中
$\overline{\mathfrak F}_{k,\mathrm{red}}$ 的定义域写成纯约化张量模
$B_k\overline{\mathfrak C}_{1,\mathrm{red}}[1]$；但条件 (4.4.31) 及证明
首步 $\overline{\mathfrak F}_{1,\mathrm{red}}(e)=e$ 都要求它定义在保留
幺元的部分约化模 $B_k^{\mathrm{red}}\overline{\mathfrak C}_1[1]$ 上。
译文依定义 4.4.22 及引理 4.4.30 统一此型别。}
\end{proposition}

\begin{proof}
证明与命题 4.4.11 类似，所以只讨论需要修改之处。归纳的第一步与引理
4.4.2 的证明相同；唯一需要核对的是
$\overline{\mathfrak F}_{1,\mathrm{red}}(e)=e$，而这自动来自它与
$\overline{\mathrm{Incl}}$ 的相容性。

假设命题 4.4.34 对 $K-1$ 已经成立。也就是说，已有
$\overline{\mathfrak F}_{k,\mathrm{red}}$（$k=1,\ldots,K-1$）。于是
得到 $\overline{\mathfrak F}^{(K-1)}$，它诱导
\[
 \widehat{\overline{\mathfrak F}}_{\mathrm{red}}^{(K-1)}\colon
 B_{1\cdots K}^{\mathrm{red}}\overline{\mathfrak C}_1[1]
 \longrightarrow B_{1\cdots K}\overline{\mathfrak C}_2[1].
\]
考虑
\[
 \begin{aligned}
 \overline{\operatorname{Err}}_K
 &={}
 \widehat{\overline d}\circ
 \widehat{\overline{\mathfrak F}}_{\mathrm{red}}^{(K-1)}
 -\widehat{\overline{\mathfrak F}}_{\mathrm{red}}^{(K-1)}
 \circ\widehat{\overline d}^{\mathrm{red}}\\
 &\in\operatorname{Hom}\bigl(
 B_{1\cdots K}^{\mathrm{red}}\overline{\mathfrak C}_1[1],
 B_{1\cdots K}\overline{\mathfrak C}_2[1]
 \bigr).
 \end{aligned}
\]
与命题 4.4.11 的证明相同，可以证明
\[
 \overline{\operatorname{Err}}_K
 \in\operatorname{Hom}\bigl(
 B_K^{\mathrm{red}}\overline{\mathfrak C}_1[1],
 \overline{\mathfrak C}_2[1]
 \bigr),
\]
\[
 \begin{aligned}
 \overline{\operatorname{Err}}_K
 &={}
 \overline{\mathfrak m}\circ
 \widehat{\overline{\mathfrak F}}_{\mathrm{red}}^{(K-1)}
 -\overline{\mathfrak F}_{\mathrm{red}}^{(K-1)}
 \circ\widehat{\overline d}^{\mathrm{red}}\\
 &\in\operatorname{Hom}\bigl(
 B_{1\cdots K}^{\mathrm{red}}\overline{\mathfrak C}_1[1],
 B_{1\cdots K}\overline{\mathfrak C}_2[1]
 \bigr),
 \end{aligned}
\]
并且
\[
 \delta_1(\overline{\operatorname{Err}}_K)=0,
 \qquad
 \overline{\operatorname{Err}}_K\in\operatorname{Im}\delta_1.
\]
这里，$\overline{\mathfrak F}_{\mathrm{red}}^{(K-1)}$ 是从
$B_{1\cdots K}^{\mathrm{red}}\overline{\mathfrak C}_1[1]$ 到
$\overline{\mathfrak C}_2[1]$ 的同态；它在
$B_k^{\mathrm{red}}\overline{\mathfrak C}_1[1]$ 上等于
$\overline{\mathfrak F}_{k,\mathrm{red}}$（$k=1,\ldots,K-1$），而在
$B_K^{\mathrm{red}}\overline{\mathfrak C}_1[1]$ 上为零。
\translatornote{原书把括号内的分量范围写成 $k=1,\ldots,K$，但紧接着
又规定第 $K$ 个分量为零；归纳假设也只给出到 $K-1$ 为止的分量，故译文
改为 $k=1,\ldots,K-1$。}
需要新核对的是下面一点。

\begin{lemma}{4.4.35}\label{lem:4-4-35}
\[
 \overline{\operatorname{Err}}_K\circ s_{\mathrm{left}}
 =\overline{\operatorname{Err}}_K\circ s_{\mathrm{right}}=0.
\]
\end{lemma}

\begin{proof}
先计算
\begin{equation}
 \begin{aligned}
 \bigl(
 \widehat{\overline{\mathfrak F}}_{\mathrm{red}}^{(K-1)}
 \circ s_{\mathrm{right}}\bigr)(\mathbf x)
 &={}
 (-1)^{\deg\mathbf x}
 \widehat{\overline{\mathfrak F}}_{\mathrm{red}}^{(K-1)}(\mathbf x)
 \otimes\overline{\mathfrak F}_{1,\mathrm{red}}(e)\\
 &=\bigl(s_{\mathrm{right}}\circ
 \widehat{\overline{\mathfrak F}}_{\mathrm{red}}^{(K-1)}
 \bigr)(\mathbf x),
 \end{aligned}
\tag{4.4.36}\label{eq:4-4-36}
\end{equation}
其中用了 \eqref{eq:4-4-31-2}。因而由引理 4.4.25 和
\eqref{eq:4-4-36}，有
\begin{equation}
 \begin{aligned}
 \overline{\operatorname{Err}}_K\circ s_{\mathrm{right}}
 &={}
 \overline{\mathfrak m}\circ
 \widehat{\overline{\mathfrak F}}_{\mathrm{red}}^{(K-1)}
 \circ s_{\mathrm{right}}
 -\overline{\mathfrak F}_{\mathrm{red}}^{(K-1)}
 \circ\widehat{\overline d}^{\mathrm{red}}
 \circ s_{\mathrm{right}}\\
 &=\overline{\mathfrak m}\circ s_{\mathrm{right}}\circ
 \widehat{\overline{\mathfrak F}}_{\mathrm{red}}^{(K-1)}
 +\overline{\mathfrak F}_{\mathrm{red}}^{(K-1)}
 \circ s_{\mathrm{right}}\circ
 \widehat{\overline d}^{\mathrm{red}}
 -\overline{\mathfrak F}_{\mathrm{red}}^{(K-1)}.
 \end{aligned}
\tag{4.4.37}\label{eq:4-4-37}
\end{equation}
由 \eqref{eq:4-4-31-2}，有
$\overline{\mathfrak F}_{\mathrm{red}}^{(K-1)}
\circ s_{\mathrm{right}}=0$。（注意这个公式中的
$\overline{\mathfrak F}_{\mathrm{red}}^{(K-1)}$ 没有帽号。）另一方面，
由 $s_{\mathrm{right}}$ 的定义，
\begin{equation}
 \begin{aligned}
 &\bigl(
 \overline{\mathfrak m}\circ s_{\mathrm{right}}\circ
 \widehat{\overline{\mathfrak F}}_{\mathrm{red}}^{(K-1)}
 \bigr)(\mathbf x)\\
 &\qquad={}
 (-1)^{\deg\mathbf x}\overline{\mathfrak m}
 \bigl(
 \widehat{\overline{\mathfrak F}}_{\mathrm{red}}^{(K-1)}(\mathbf x)
 \otimes e
 \bigr).
 \end{aligned}
\tag{4.4.38}\label{eq:4-4-38}
\end{equation}
右端只有
$B_1\overline{\mathfrak C}_2[1]\otimes e$ 中的项非零，所以
\eqref{eq:4-4-38} 等于
$\overline{\mathfrak F}_{\mathrm{red}}^{(K-1)}(\mathbf x)$。因此
\eqref{eq:4-4-37} 为零。等式
$\overline{\operatorname{Err}}_K\circ s_{\mathrm{left}}=0$ 的证明类似。
\end{proof}

其余证明与命题 4.4.11 的证明类似。利用引理 4.4.35，可以把
\[
 \overline{\operatorname{Err}}_K
 \in\operatorname{Hom}\bigl(
 B_K\overline{\mathfrak C}_{1,\mathrm{red}}[1],
 \overline{\mathfrak C}_2[1]
 \bigr)
\]
看作一个同态。可以用 $\overline{\mathfrak m}_1$ 在
$\operatorname{Hom}(B_K\overline{\mathfrak C}_{1,\mathrm{red}}[1],
\overline{\mathfrak C}_2[1])$ 上定义上边界算子 $\delta_1$。（不能用
$\overline{\mathfrak m}_2$ 等在
$B\overline{\mathfrak C}_{\mathrm{red}}[1]$ 上定义算子。）

于是可以照搬命题 4.4.11 证明中的论证，得到
$\overline{\operatorname{Err}}_K\in\operatorname{Im}\delta_1$。因而存在
\[
 \overline{\mathfrak F}'_{K,\mathrm{red}}
 \in\operatorname{Hom}\bigl(
 B_K\overline{\mathfrak C}_{1,\mathrm{red}}[1],
 \overline{\mathfrak C}_2[1]
 \bigr),
 \qquad
 \delta_1\overline{\mathfrak F}'_{K,\mathrm{red}}
 +\overline{\operatorname{Err}}_K=0.
\]
它可以视为满足 \eqref{eq:4-4-31-2} 的同态
\translatornote{原书先写
$\overline{\operatorname{Err}}_K
=\delta_1\overline{\mathfrak F}'_{K,\mathrm{red}}$，但本段沿用命题
4.4.11 的误差项定义，扩张条件应为
$\delta_1\overline{\mathfrak F}'_{K,\mathrm{red}}
+\overline{\operatorname{Err}}_K=0$；原书紧接着又把上述同态的靶误写成
$\overline{\mathfrak C}_1[1]$，而前一句的同态空间及整个构造都唯一要求
靶为 $\overline{\mathfrak C}_2[1]$。译文据此一并修正。}
\[
 B_K^{\mathrm{red}}\overline{\mathfrak C}_1[1]
 \longrightarrow\overline{\mathfrak C}_2[1].
\]
利用它可得到含幺 $A_K$ 同态
$\overline{\mathfrak F}^{(K)}\colon
\overline{\mathfrak C}_1\to\overline{\mathfrak C}_2$。再完全照命题
4.4.11 的证明修改它，使它与 $\overline{\mathrm{Incl}}$、
$\overline{\mathrm{Eval}}$ 相容。命题 4.4.34 得证。
\end{proof}

命题 4.4.34 立即蕴含定理 4.2.34 的无滤过含幺版本。
